\section{Project Scope}

The aim of this project is not to announce a decipherment prematurely. The aim is to build the strongest possible evidence pipeline: corpus hygiene, sign statistics, positional grammar, non-linguistic baselines, and only then constrained linguistic hypotheses.

\subsection{Working Stance}

\begin{itemize}
  \item Treat the Indus system as undeciphered.
  \item Treat all proposed readings as hypotheses, not data.
  \item Prefer reproducible tests over attractive resemblances.
  \item Compare Dravidian, Indo-Aryan/Indo-Iranian, Munda-like, isolate, and non-linguistic models under the same evidentiary rules.
  \item Separate graphemic questions from linguistic questions: first ask what the signs are doing, then ask whether speech is involved.
\end{itemize}

\subsection{Current Data}

\begin{itemize}
  \item \texttt{data/ivs\_corpus\_cleaned.csv}: initial neutral working corpus.
  \item \texttt{data/ivs\_corpus.csv}: same row count, but includes proposed Sanskrit readings and translations. These fields are excluded from the baseline because they risk circular reasoning.
\end{itemize}

\subsection{Research Phases}

\begin{enumerate}
  \item Corpus audit and provenance.
  \item Sign inventory normalization.
  \item Directionality and positional class analysis.
  \item Formula, n-gram, and network analysis.
  \item Non-linguistic baseline modeling.
  \item Linguistic hypothesis testing.
  \item Narrow, falsifiable candidate readings only where distributional evidence permits them.
\end{enumerate}

\subsection{Evidence Ladder}

\begin{description}
  \item[Speculative] Plausible idea, not yet tested.
  \item[Suggestive] Agrees with some observed pattern but has weak alternatives.
  \item[Testable] Predicts a measurable corpus pattern.
  \item[Statistically supported] Survives predefined quantitative tests.
  \item[Philologically plausible] Agrees with linguistic, archaeological, and comparative constraints.
\end{description}

No claim enters the project as a conclusion. It enters as a hypothesis and earns its place.

\begin{figure}[htbp]
\centering
\begin{tikzpicture}[
  node distance=0.9cm and 1.1cm,
  stage/.style={
    rectangle,
    rounded corners=2pt,
    draw=blue!45!black,
    fill=blue!4,
    text width=3.0cm,
    align=center,
    minimum height=0.9cm,
    font=\small
  },
  arrow/.style={-{Latex[length=2.5mm]}, thick, draw=blue!45!black}
]
\node[stage] (corpus) {Corpus audit};
\node[stage, right=of corpus] (inventory) {Sign inventory};
\node[stage, right=of inventory] (grammar) {Positional grammar};
\node[stage, below=of grammar] (baseline) {Rival baselines};
\node[stage, left=of baseline] (linguistic) {Linguistic tests};
\node[stage, left=of linguistic] (readings) {Candidate readings};
\draw[arrow] (corpus) -- (inventory);
\draw[arrow] (inventory) -- (grammar);
\draw[arrow] (grammar) -- (baseline);
\draw[arrow] (baseline) -- (linguistic);
\draw[arrow] (linguistic) -- (readings);
\end{tikzpicture}
\caption{Research pipeline. Candidate readings are deliberately last.}
\label{fig:pipeline}
\end{figure}

